语句集合在某种意义上描述了它们在其中共同为真的那些结构；这些结构就是它们的\emph{模型}。
反过来，如果我们从一个结构或一组结构出发，可能会对以它们为模型的语句集合感兴趣，这就是该结构或结构组的\emph{理论}。
任何这样的语句集合都具有这样的性质：每一个被它们语义地后承的语句都已经在这个集合中；它们是\emph{封闭的}。
更一般地，如果一个集合 $\Gamma$ 在语义后承下封闭，我们就称 $\Gamma$ 为一个理论；如果 $\Gamma$ 由被 $\Delta$ 语义地后承的所有语句组成，则称 $\Gamma$ 由~$\Delta$ \emph{公理化}。

数学提供了许多理论的例子，例如线性序理论、群理论，或算术理论（如由 Peano（皮亚诺）公理公理化的理论）。
但在其他学科中也有许多重要理论的例子，例如，关系数据库可以被视为理论，形而上学也关注可公理化的部分论理论。

当我们着手研究一个理论时，一个重要的问题是它的语言是否有足够的表达力，能让我们表述所有希望该理论讨论的内容；
另一个问题是它是否足够强大，能证明我们希望它证明的结论。
要\emph{表达}一个关系，我们需要一个具有所需数量自由变元的公式。在\emph{集合论}中，我们只有 $\in$ 作为关系符号，
但它允许我们用 $\lforall[u][(u \in x \lif u \in y)]$ 来表达 $x \subseteq y$。事实上，\emph{Zermelo（策梅洛）-Fraenkel（弗兰克尔）集合论}~$\Log{ZFC}$ 仅凭少数几条公理和一系列日益复杂的定义（如 $\subseteq$ 的定义）构成的链条，就足够强大，既能（几乎）表达每一个数学命题，又能（几乎）证明每一个数学定理。